% =====================================================================
%  Measurement Validity is not Decision Value
%  A validated market regime instrument and the honest limits of its use
%
%  COLLABORATIVE DRAFT.  Read README.md for the editing conventions.
%  - Blue text = AI-drafted, NOT yet reviewed by Adam. Wrapped in \aidraft{}.
%    As you rewrite a passage in your own voice, delete the \aidraft{...}
%    wrapper (keep the text) and it turns black = "approved".
%  - Flip \reviewtrue -> \reviewfalse below for a clean reading/submission PDF.
% =====================================================================
\documentclass[11pt]{article}

% ---- REVIEW TOGGLE ---------------------------------------------------
% \reviewtrue  : show AI-draft highlighting, margin TODOs, and a checklist page
% \reviewfalse : clean render (no colors, no todos) for reading / submission
\newif\ifreview
\reviewtrue
% ----------------------------------------------------------------------

\usepackage[margin=1in]{geometry}
\usepackage[T1]{fontenc}
\usepackage{lmodern}
\usepackage{microtype}
\usepackage{amsmath,amssymb}
\usepackage{booktabs}
\usepackage{xcolor}
\usepackage{csquotes}
\usepackage{graphicx}
\usepackage[authoryear,round]{natbib}
\usepackage{hyperref}
\hypersetup{colorlinks=true,linkcolor=black,citecolor=blue!45!black,urlcolor=blue!45!black}

% ---- review tooling --------------------------------------------------
\definecolor{aidraftcolor}{HTML}{1A56A0}   % AI-draft prose color (blue)
\ifreview
  \usepackage[colorinlistoftodos,prependcaption,textsize=footnotesize]{todonotes}
  \newcommand{\aidraft}[1]{{\color{aidraftcolor}#1}} % unreviewed AI prose
  \newcommand{\needcite}[1]{\todo[color=orange!55,size=\tiny]{CITE: #1}}
  \newcommand{\reviewnote}[1]{\todo[color=yellow!75,size=\tiny]{#1}}
  \newcommand{\stub}[1]{\todo[inline,color=red!25]{STUB / TO WRITE: #1}}
\else
  \usepackage[disable]{todonotes}
  \newcommand{\aidraft}[1]{#1}
  \newcommand{\needcite}[1]{}
  \newcommand{\reviewnote}[1]{}
  \newcommand{\stub}[1]{}
\fi

% ---- single-sourced frozen numbers (edit once, used everywhere) ------
% >>> begin numbers.tex
% =====================================================================
%  numbers.tex --- single source of truth for every frozen figure.
%  Edit a number HERE and it updates everywhere it is used.
%  Sources: results/stage1.csv, allocation_summary.csv, calibration_gate.csv,
%           sensor_validation.csv, rotation_* (git d022c09), trend_proxy_gate.csv.
%  FKO = Fleming-Kirby-Ostdiek quadratic-utility fee, annualized bps/yr,
%  90% CIs from paired stationary bootstrap (mean block 126d, B=2000, seed 0).
% =====================================================================

% ---- Instrument (Layer 1) ----
\newcommand{\stabJM}{1.000}          % label stability, +/-2y train shift
\newcommand{\stabHMM}{0.809}         % HMM ensemble self-agreement
\newcommand{\switchesYr}{1.66}
\newcommand{\nEpisodes}{30}
\newcommand{\calmVol}{14\%}
\newcommand{\stressVol}{25\%}
\newcommand{\bearsCaught}{15 of 18}
\newcommand{\detLag}{20}              % median recognition lag, trading days
\newcommand{\spliceCorr}{0.9957}
\newcommand{\overlapAgree}{1.0000}
\newcommand{\overlapDays}{9{,}120}

% ---- Race 1: exposure overlay ----
\newcommand{\feeJMBH}{+487.8}
\newcommand{\feeJMBHci}{[-47.4,\,+1138.8]}
\newcommand{\feeJMVT}{-255.8}
\newcommand{\feeJMVTci}{[-477.5,\,-36.7]}
\newcommand{\feeJMBHgone}{-314.2}    % gamma=1
\newcommand{\feeStaticMix}{-48.7}
\newcommand{\placeboOne}{541.3}      % C1 95th pctile
\newcommand{\breakeven}{305.6}

% ---- Race 2: conditional covariance ----
\newcommand{\feeAmatch}{+2.8}
\newcommand{\feeAmatchci}{[-3.2,\,+8.4]}
\newcommand{\feeBreact}{-29.1}
\newcommand{\feeBreactci}{[-66.8,\,+8.5]}
\newcommand{\volofvolCond}{2.11}
\newcommand{\volofvolMatch}{2.15}

% ---- Race 3: confidence margin ----
\newcommand{\dgpCells}{eight of eight}

% ---- Race 4: cross-asset rotation (Path-B null, descriptive) ----
\newcommand{\trendProxyCorr}{0.69}   % TSMOM proxy vs DBMF
\newcommand{\sharpeTimed}{0.62}
\newcommand{\sharpeStatic}{0.70}
\newcommand{\sharpeVT}{0.66}          % VT matched-exposure
\newcommand{\mixedDays}{76\%}         % per-asset regimes in mixed states

% <<< end numbers.tex


% ---- metadata --------------------------------------------------------
\title{Measurement Validity is not Decision Value:\\
       A Validated Market Regime Instrument and the Honest Limits of Its Use}
\author{Adam Morris\thanks{Working draft. Not for circulation.}}
\date{Draft --- \today}

\begin{document}
\maketitle

\ifreview
  \begin{center}\small\itshape
  Review render: blue text is AI-drafted and awaiting your edit; margin notes are action items.
  Flip \texttt{\textbackslash reviewtrue} to \texttt{\textbackslash reviewfalse} in main.tex for a clean PDF.
  \end{center}
  \listoftodos
  \clearpage
\fi

% >>> begin sections/00-abstract.tex
\begin{abstract}
\aidraft{%
We build and validate a market regime-detection instrument --- a two-state statistical jump
model on return-only downside features --- to an unusual standard, and then map, under strict
preregistered one-look discipline, exactly where its information is and is not usable for
decisions.
As a measurement device the label is exceptional: \stabJM{} agreement under $\pm 2$-year
training-window perturbation (versus \stabHMM{} for an HMM ensemble), \switchesYr{} switches
per year, and live operation to the present.
As a decision input it is dominated at every use we tested by a reactive estimator that never
names the regime: binary exposure timing loses \feeJMVT{}~bps/yr to volatility targeting;
covariance conditioning loses \feeBreact{}~bps/yr to a plain EWMA covariance; the filter's own
confidence margin loses to an EWMA-volatility signal in \dgpCells{} synthetic environments; and
cross-asset defensive rotation adds no risk-adjusted value beyond exposure, its drawdown
protection matched by volatility targeting.
A single mechanism, shown in simulation before any real-data run, unifies the nulls: regime
value at daily horizons is an information-timing claim, and a causal detector's recognition lag
lands it precisely where even a perfectly informed oracle loses to an estimator that pays no
lag toll.
The instrument-versus-strategy distinction, usually a footnote, is the thesis: measurement
validity is not decision value, and the applied regime-switching literature's positive claims
live in the gap between them.%
}
\end{abstract}

% <<< end sections/00-abstract.tex

% >>> begin sections/01-intro.tex
\section{Introduction}
\label{sec:intro}

\aidraft{%
A recurring claim in the applied regime-switching literature is that identifying the market's
latent state and trading on it improves investor outcomes.
Recent work built on statistical jump models \citep{nystrup2020eswa,nystrup2021,shu2024,jmmpc2025}
reports that regime-aware allocations beat buy-and-hold and raise risk-adjusted returns.
The claim is attractive because the estimator is attractive: jump models produce stable,
interpretable state sequences from a handful of return-based features, without the fragility
that made hidden Markov regime models hard to deploy out of sample.%
}
\needcite{Shu--Yu--Mulvey exact headline improvement + their benchmark}
\aidraft{%
The cleanest instance is the allocation sequel of \citet{shu2025dynamic}: regime forecasts enter
expected returns through a gradient-boosted classifier, the covariance is a plain EWMA, and the
benchmarks are minimum-variance, mean-variance, and equal-weight --- with no exposure-matched,
volatility-targeted, or cost-matched control.
\reviewnote{Adam: confirm the 60/20/20-vs-60/40 framing you want here.}%
}

\aidraft{%
We take the estimator seriously as a measurement device and separate two questions the
literature tends to fuse.
First: does the instrument \emph{measure} the market's risk state reliably --- is the label a
stable, reproducible, live-computable quantity?
Second: does \emph{conditioning a decision} on that label add value over a simple estimator that
never names the regime?
Our central finding is that these questions have opposite answers, and that the reported
economic value lives in the gap between them.%
}

\aidraft{%
On the first question the instrument is exceptional.
A $K{=}2$ jump model on return-only downside features produces a label that is \stabJM{} stable
under $\pm 2$-year perturbations of the training window --- where an HMM ensemble on the same
data agrees with itself only \stabHMM{} of the time --- switches a modest \switchesYr{} times per
year, and reproduces to \overlapAgree{} agreement on a \overlapDays{}-day out-of-sample overlap
when carried to the present through a live data splice.
As a market-state sensor it is genuine: it flags \bearsCaught{} ex-post bear markets.
This is a measurement instrument worth having, and Section~\ref{sec:instrument} validates it as such.%
}

\aidraft{%
On the second question every use we tested is dominated by a reactive estimator computing the
decision-relevant quantity directly.
We evaluate four uses under strict preregistered, one-look discipline (Section~\ref{sec:discipline}):
trading the label as an exposure signal (Section~\ref{sec:race1}), conditioning a multi-asset
covariance estimate on it (Section~\ref{sec:race2}), mapping the filter's own evidence margin to
a state probability (Section~\ref{sec:race3}), and rotating a defensive cross-asset sleeve on it
(Section~\ref{sec:race4}).
The exposure strategy's apparent \feeJMBH{}~bps/yr utility fee over buy-and-hold sits inside
every exposure-matched null band and loses a significant \feeJMVT{}~bps/yr to volatility
targeting; covariance conditioning adds a band-interior \feeAmatch{}~bps/yr and loses
\feeBreact{}~bps/yr to a plain EWMA covariance; the evidence margin is beaten by an EWMA-volatility
probability in \dgpCells{} synthetic environments; and cross-asset rotation's drawdown protection
is matched by a volatility target at equal average exposure.
A single mechanism, demonstrated in simulation before any real-data run, unifies them
(Section~\ref{sec:mechanism}): regime value at daily horizons is an information-\emph{timing}
claim, and a causal detector's recognition lag lands it exactly where even a perfectly informed
oracle loses to an estimator that pays no recognition toll.%
}

\aidraft{%
Our contribution is therefore both constructive and deflationary.
We deliver a regime instrument validated to an unusual standard --- perturbation stability, live
operation, and a claim-by-claim ledger of what is out-of-sample-validated versus merely
descriptive (Section~\ref{sec:monitor}, the monitor).
And we show, with exposure-matched nulls and reactive co-primary baselines that the source
literature omits, that its decision value at daily horizons is an artifact of comparison choices.
The instrument-versus-strategy distinction --- usually a footnote --- is the paper's thesis:
measurement validity is not decision value, and honest baselines separate them.%
}

% <<< end sections/01-intro.tex

% >>> begin sections/02-discipline.tex
\section{Discipline}
\label{sec:discipline}

\aidraft{%
Every economic claim in this paper was evaluated exactly once, against criteria frozen before any
real-data result existed.
We preregister each experiment in a versioned document (objects, arms, primary metric, falsifiers,
and verdict wording), freeze it in the project's public git history, and only then run: the freeze
commit provably precedes the one-look run in the commit graph.
Results cannot be re-run, re-classified, or re-worded after the look; where a design defect is
discovered post hoc --- it happened once, a qualitatively-specified gate criterion --- it is
recorded as a defect rather than resolved silently.
Two further rules harden the protocol beyond common practice.
First, a cooling-off requirement: any experiment that could produce a positive claim requires an
overnight gap between freeze and run, and an explicit signed authorization --- never a
conversational go-ahead.
Second, an out-of-hypothesis-sample rule: hypotheses generated by descriptive work on the U.S.\
panel are treated as contaminated by selection, and no supportive verdict may be claimed for them
until they confirm on data that did not generate them (designated developed-market panels, held
untouched).%
}

\aidraft{%
The estimation architecture is causal end to end.
Features are derived from the return series alone --- exponentially weighted downside deviation
(10-day half-life) and Sortino ratios (20- and 60-day half-lives) --- eliminating the
vintage-lookahead problem class that macro conditioning series import.
The jump model is refit annually on an expanding window; within each scoring year every parameter
is frozen, including the jump penalty $\lambda$, which is re-chosen each refit by causal validation
on the trailing eight years of the training window only.
Out-of-sample states come from a forward dynamic-programming filter --- the endpoint of the value
recursion, never the backtracked (smoothed) path --- so the state at time $t$ depends only on data
through $t$.
Execution is next-close (signals trade with a two-day shift; same-close is reported as sensitivity
only), with 10~bps one-way costs identical across every arm.%
}

\aidraft{%
Inference is comparative and exposure-honest.
The primary metric is the Fleming--Kirby--Ostdiek quadratic-utility performance fee
\citep{fko2001} ($\gamma{=}10$ headline, $\gamma{=}1$ sensitivity) with paired
stationary-bootstrap confidence intervals \citep{politis1994} (mean block 126 days).
Because utility fees mechanically reward average exposure, every conditional arm faces three null
distributions computed before the primary: random persistent signals with matched transition rates
(exposure-matched placebos), timing-destroyed surrogates, and iid pipelines --- plus a static mix
matched to the strategy's realized average exposure.
And because \enquote{beats buy-and-hold} is not evidence of skill, every experiment carries a
reactive incumbent as co-primary: volatility targeting for exposure decisions, EWMA estimation for
covariance decisions.
A claim counts only if it clears the best simple alternative; arms adjacent in the design differ by
exactly one ingredient, so any gap is attributable.%
}

% <<< end sections/02-discipline.tex

% >>> begin sections/03-instrument.tex
\section{The instrument (Layer 1 --- measurement validity)}
\label{sec:instrument}

\aidraft{%
The estimator is a two-state weighted statistical jump model \citep{bemporad2018,nystrup2020eswa}
fit to three return-only features --- exponentially weighted downside deviation (10-day half-life)
and Sortino ratios at 20- and 60-day half-lives --- with a jump penalty $\lambda$ that enforces
state persistence.
We evaluate it first purely as a measurement device, before any economic use, on three properties
a deployable instrument must have: stability under reasonable estimation choices, agreement with an
external ground truth, and live computability.%
}

\paragraph{Stability.}
\aidraft{%
The property that motivated the whole program is reproducibility of the label under
training-window perturbation.
Shifting the training start by $\pm 2$ years --- a first-order analyst choice that leaves competing
estimators disagreeing about the state on a fifth of all days --- changes the jump-model label on
$0.0\%$ of out-of-sample days: \stabJM{}/\stabJM{} agreement at both the $-2$y and $+2$y shifts.
An HMM ensemble on the identical data and features achieves only \stabHMM{}.
The label switches \switchesYr{} times per year and partitions the 1990--2026 out-of-sample period
into \nEpisodes{} episodes, each a genuine high-volatility environment (calm-state annualized
volatility ${\approx}\,\calmVol$, stressed-state ${\approx}\,\stressVol$); $\lambda$ selects to the
interior of its search grid rather than an edge.
Estimation instability --- the practical reason HMM regime labels are hard to trust in production
--- is, for this estimator on this data, solved.%
}
\reviewnote{Adam: 1.000 is under-argued as it stands. A constant label is also perfectly stable ---
so pair stability with skill; harden the perturbation ($\pm$5y, bootstrap-resample, param jitter)
and show the degradation curve; add a deterministic vol-threshold baseline next to the stochastic
HMM. This is the referee's first objection. (See NOTES.md instrument to-do.)}

\paragraph{External agreement, and an honest boundary.}
\aidraft{%
Validated against ex-post bear-market datings, the label catches \bearsCaught{} declines of
$-15\%$ or more, at a median recognition lag of \detLag{} trading days from the ex-post peak.
It is, precisely, a \emph{volatility-state} detector rather than a \emph{bear-market} detector: it
fires on high-volatility environments whether or not a sustained decline follows (hence a moderate
precision against directional bear datings), and it misses fast crashes that resolve before its
persistence-seeking filter can accumulate evidence --- 1998 (LTCM) and 2018-Q4 are sub-threshold at
the deployed $\lambda$.
We state this as the instrument's boundary of use rather than a defect to be tuned away: the
\detLag{}-day lag and the fast-crash misses are exactly the constraints the mechanism section shows
are binding for any causal detector, and they define what the monitor (Section~\ref{sec:monitor})
can and cannot claim.%
}

\paragraph{Live computability.}
\aidraft{%
The instrument runs to the present.
Ken French publishes with a one-to-two-month lag, so the live tail is spliced from an SPY proxy;
the splice passes its gate at \spliceCorr{} return correlation, and the composite label agrees with
the frozen research label on \overlapAgree{} of the \overlapDays{}-day overlap.
There is no vintage-lookahead: the filter's state at time $t$ uses only data through $t$, and the
live label at $t$ is the same object the frozen backtest would have produced.%
}

\aidraft{%
These three properties make the label a measurement instrument worth building on.
The remaining sections ask what happens when a decision is conditioned on it --- and find,
uniformly, that the measurement's quality does not transfer to decision value.%
}

% <<< end sections/03-instrument.tex

% >>> begin sections/04-race1-exposure.tex
\section{Race 1 --- trading the label directly}
\label{sec:race1}

\aidraft{%
The first and most direct use is the one the literature reports: hold the risky asset in the calm
state, de-risk in the stressed state.
We run the $K{=}2$ jump model on French daily total returns (trained from 1970, scored
1990-03--2026-05, $n{=}\overlapDays{}$ out-of-sample days) and evaluate the overlay against
buy-and-hold, against volatility targeting as the honest reactive incumbent, and against three
exposure-matched null distributions.%
}

\aidraft{%
The headline looks like the literature's: a \feeJMBH{}~bps/yr quadratic-utility fee over
buy-and-hold (Table~\ref{tab:exp1}).
It does not survive a single honest control.
The point estimate sits \emph{inside} every exposure-matched null band --- the persistent-placebo
95th percentile alone is $\placeboOne$, above the observed fee --- so random signals with matched
transition rates earn as much.
An exposure-matched static mix beats the overlay outright ($\feeStaticMix$), and at $\gamma{=}1$ the
fee turns to $\feeJMBHgone$.
Against the honest bar the verdict is unambiguous: $\text{fee}(\text{JM}-\text{VT}) = \feeJMVT$~bps/yr,
90\% CI $\feeJMVTci$, a significant loss to volatility targeting.
The frozen classification is Case~B: the beat-buy-and-hold result is an exposure/utility artifact,
and timing loses to the reactive incumbent.%
}

% >>> begin tables/tab-exp1.tex
\begin{table}[t]
\centering
\small
\caption{Race 1 --- exposure overlay on equity: primary and controls. FKO quadratic-utility fee,
annualized bps/yr, $\gamma{=}10$ unless noted; 90\% CIs from the paired stationary bootstrap
(mean block 126d, $B{=}2000$). Source: \texttt{results/stage1.csv}.}
\label{tab:exp1}
\begin{tabular}{lr}
\toprule
Quantity & Value \\
\midrule
fee(JM $-$ B\&H), $\gamma{=}10$              & \feeJMBH~\feeJMBHci \\
fee(JM $-$ VT), $\gamma{=}10$ (honest bar)   & \feeJMVT~\feeJMVTci \\
fee(JM $-$ B\&H), $\gamma{=}1$               & \feeJMBHgone \\
fee(JM $-$ exposure-matched static mix)      & \feeStaticMix \\
same-close (delay $=1$) sensitivity          & $+480.1$ \\
break-even one-way cost                       & \breakeven~bps \\
C1 placebo 95th percentile ($n{=}100$)        & $+\placeboOne$ \\
C2 surrogate band max ($n{=}10$)              & $+1080.0$ \\
C3 iid band max ($n{=}10$)                    & $+805.3$ \\
label stability, train-start $+2$y / $-2$y    & \stabJM{} / \stabJM{} \\
era-half fees 1990--2007 / 2008--2026         & $+237.9$ / $+729.5$ \\
switches/yr                                    & \switchesYr{} \\
frozen classification                         & Case B \\
\bottomrule
\end{tabular}
\end{table}

% <<< end tables/tab-exp1.tex


% <<< end sections/04-race1-exposure.tex

% >>> begin sections/05-race2-covariance.tex
\section{Race 2 --- conditioning covariance on the label}
\label{sec:race2}

\aidraft{%
If timing exposure fails, perhaps the state helps estimate \emph{risk} rather than direction.
We build a long-only equal-risk-contribution portfolio over equity, ten-year Treasury total return,
gold, and cash, targeting 8\% volatility (scale-down only), and condition its covariance estimate on
the frozen state label (hard state, expanding per-state windows, 0.5 shrinkage to unconditional),
rebalancing monthly and on state flips.
The matched control is the identical allocator with conditioning removed; the reactive co-primary is
the same allocator on a plain EWMA covariance ($\lambda{=}0.97$).%
}

\aidraft{%
Conditioning adds essentially nothing over its matched twin --- $\feeAmatch$~bps/yr, CI
$\feeAmatchci$, inside all null bands (Table~\ref{tab:exp2}) --- and loses to the reactive
incumbent: $\text{fee}_B(\text{cond}-\text{EWMA}) = \feeBreact$~bps/yr, with the EWMA arm the best
performer outright (Table~\ref{tab:exp2arms}).
The mechanism the state is supposed to exploit is real but second-best: conditioning does stabilize
portfolio risk (vol-of-vol $\volofvolCond$ vs $\volofvolMatch$~pp for the matched twin), it is simply
dominated at the job by an estimator that re-estimates the whole covariance in weeks while the
expanding per-state estimate drags decades of stale mixture.
The frozen classification is null.%
}

% >>> begin tables/tab-exp2.tex
\begin{table}[t]
\centering
\small
\caption{Race 2 --- conditional allocation: primary and controls. FKO fee, annualized bps/yr,
$\gamma{=}10$. Source: \texttt{results/allocation\_summary.csv}, \texttt{allocation\_controls.csv}.}
\label{tab:exp2}
\begin{tabular}{lr}
\toprule
Quantity & Value \\
\midrule
fee$_A$(cond $-$ B\_match), full window       & \feeAmatch~\feeAmatchci \\
fee$_B$(cond $-$ EWMA), from activation       & \feeBreact~\feeBreactci \\
fee$_A$ / fee$_B$ at $\gamma{=}1$             & $+1.2$ / $+3.2$ \\
placebo ($n{=}100$): fee$_A$ 5--95th          & $[-0.5,\,+5.8]$ \\
placebo ($n{=}100$): fee$_B$ 5--95th          & $[-32.7,\,-25.7]$ \\
episode-shuffle ($n{=}50$): fee$_A$ 5--95th   & $[-1.7,\,+4.9]$ \\
iid panels ($n{=}20$): fee$_A$ 5--95th        & $[-1.2,\,+0.9]$ \\
vol-of-vol, cond / B\_match (F2)              & \volofvolCond{} / \volofvolMatch{}~pp \\
turnover/yr, cond / B\_match / EWMA          & 0.11 / 0.06 / 0.88 \\
conditioning activation                       & 1994-03-29 \\
falsifiers F1 / F1b / F2 / F3                  & T / T / F / F \\
frozen classification                         & NULL \\
\bottomrule
\end{tabular}
\end{table}

% <<< end tables/tab-exp2.tex

% >>> begin tables/tab-exp2-arms.tex
\begin{table}[t]
\centering
\small
\caption{Race 2 --- arm summary over the scored window. The reactive EWMA arm is the best performer;
buy-and-hold equity is shown for exposure context. Source: \texttt{results/allocation\_arms.csv}.}
\label{tab:exp2arms}
\begin{tabular}{lrrrr}
\toprule
Arm & Ann.\ return & Ann.\ vol & Sharpe & MaxDD \\
\midrule
Conditional ERC              & 8.34\% & 6.88\%  & 0.826 & $-18.3\%$ \\
B\_match (uncond.\ expanding) & 8.33\% & 6.90\%  & 0.822 & $-18.3\%$ \\
B\_react (EWMA $\lambda{=}0.97$) & 8.28\% & 6.40\% & \textbf{0.878} & $-16.2\%$ \\
60/40                        & 9.44\% & 10.68\% & 0.634 & $-30.8\%$ \\
VT equity                    & 8.37\% & 10.04\% & 0.568 & $-26.4\%$ \\
B\&H equity                  & 12.30\% & 18.10\% & 0.533 & $-54.6\%$ \\
\bottomrule
\end{tabular}
\end{table}

% <<< end tables/tab-exp2-arms.tex


% <<< end sections/05-race2-covariance.tex

% >>> begin sections/06-race3-margin.tex
\section{Race 3 --- the confidence margin as a probability}
\label{sec:race3}

\aidraft{%
The filter carries its own confidence: the $\lambda$-normalized gap between the value of the calm and
stressed states is a natural state-probability signal.
If that margin is a portable measure of conviction, it should calibrate.
We test it on an eight-cell synthetic grid (frequency $\times$ persistence $\times$ severity $\times$
signal-to-noise), fitting a pooled Platt layer on training seeds and scoring per-cell calibration
(Cox slope and intercept), discrimination (AUC), and Brier score against two benchmarks: climatology
and a plain log-EWMA-volatility signal.%
}

\aidraft{%
The margin is not a portable confidence measure.
In \dgpCells{} data-generating environments the filter's evidence margin loses to a plain
log-EWMA-volatility signal on Brier score (by a factor of two to four) and on AUC; per-cell
calibration slopes scatter across $-0.17$ to $2.05$.
The classification is robust --- the state assignment itself is reliable --- but its internal margin
does not transfer as a probability, so the honest decision input is the hard label, not a graded
confidence.
The probability layer is killed as a decision input.%
}
\reviewnote{Adam: consider a calibration/AUC-across-cells figure here (OUTLINE figure 5).}

% <<< end sections/06-race3-margin.tex

% >>> begin sections/07-race4-crossasset.tex
\section{Race 4 --- cross-asset defensive rotation}
\label{sec:race4}

\aidraft{%
The three preceding races all time a single risk axis, where a volatility target is the obvious
reactive incumbent.
A natural objection is that the state's value is cross-sectional: even if it cannot time one asset's
exposure, per-asset regimes might diverge enough to \emph{rotate} a defensive sleeve --- a quantity no
single-axis volatility estimator can produce.
This is the design of the recent allocation literature \citep{shu2025dynamic}, and it is the one
Layer-3 use whose reactive incumbent is not, a priori, obvious.
We test it descriptively (no preregistered look was spent; the result closed the direction before one
was warranted).%
}

\aidraft{%
The cross-section exists but the timing does not add value.
Fitting a separate label per asset, the states genuinely diverge --- assets sit in mixed states on
${\approx}\,\mixedDays$ of days --- so there is something to rotate within.
Yet against a static blend that simply \emph{holds} the same diversified sleeve (equity, gold, and a
time-series-momentum trend proxy validated to \trendProxyCorr{} correlation with a live
managed-futures fund), regime-timed rotation adds no risk-adjusted value at matched average exposure:
timed Sharpe \sharpeTimed{} versus static \sharpeStatic{}, with the return difference slightly
negative after costs.
Isolating the distinctive cross-asset switch (holding equity exposure constant) leaves the static
result unchanged.
The one effect the timing does produce --- roughly halving drawdown --- is the exposure-timing
mechanism of Race~1, and it too fails the honest control: a volatility target at matched average
exposure equals or beats the timed strategy on Sharpe, drawdown, and crisis-period protection alike
(matched-exposure VT Sharpe \sharpeVT{}).
Cross-sectional selection changes the quantity being estimated, but not the verdict.%
}
\reviewnote{Adam: this is the fresh section --- the numbers are descriptive/in-sample (2007--2026,
one path). Frame it as a fourth dominated use, not a preregistered chapter. Scripts removed but
reproducible at git d022c09; findings in RESEARCH-RECORD 2026-07-26.}

% <<< end sections/07-race4-crossasset.tex

% >>> begin sections/08-mechanism.tex
\section{Mechanism --- the detection-lag race}
\label{sec:mechanism}

\aidraft{%
The results are unified by a single mechanism, demonstrated in simulation before any real-data run:
regime value at daily horizons is an information-timing claim, and recognition lag is the binding
constraint.
On synthetic two-state panels with known regimes, an oracle that switches exposure at the true
regime boundaries beats volatility targeting by $+136$ to $+1{,}090$~bps per year of utility fee ---
persistent-regime information is genuinely valuable when it is instantaneous.
Delaying the same oracle's signal by 10--21 trading days erases most to all of that edge.
Any causal detector of persistent states must accumulate evidence before switching --- our filter's
median recognition lag is on the order of 5--20 days by construction, and its empirical median
detection lag from ex-post market peaks is \detLag{} days --- so a causal regime signal lives
precisely in the regime where even perfect information loses.
The reactive incumbent has no recognition step to lose: volatility targeting responds to squared
returns within days, without ever naming the regime.
It harvests the same persistent-volatility structure the regime model detects, but pays no lag toll.%
}

\aidraft{%
The race generalizes beyond exposure timing, which is what makes it a mechanism rather than an
anecdote.
Conditioning a multi-asset covariance estimate on the state loses to a plain EWMA covariance of the
same data ($\feeBreact$~bps; the EWMA arm is the best-performing arm outright).
The model's own confidence measure loses the race: mapped through a calibration layer on synthetic
panels where the true states are volatility states, the filter's evidence margin is dominated by a
plain EWMA volatility estimate as a state-probability signal in \dgpCells{} data-generating
environments, on both discrimination and calibration.
And cross-sectional rotation, which escapes the single-axis volatility target by construction, still
sees its only real effect --- drawdown protection --- matched by that same target at equal exposure.
The pattern is uniform: the state machine knows \emph{where} it is with exceptional reliability, but
everything it knows about \emph{magnitude} --- how much risk, how confident, how to size --- arrives
later and coarser than a reactive estimator computing the same quantity directly.
Regime detection survives as measurement; regime conditioning fails as decision machinery wherever a
reactive estimator can compute the decision-relevant quantity without waiting for a label.%
}

% <<< end sections/08-mechanism.tex

% >>> begin sections/09-monitor.tex
\section{The monitor --- what the instrument is for}
\label{sec:monitor}

\stub{The constructive payoff. Draft from MONITOR-VALIDATION-SPEC.md once the monitor gate is built.}

\aidraft{%
The instrument's proper use is not a trade but a measurement: a live regime nowcast, reported with a
claim-by-claim badge of what is out-of-sample-validated versus merely descriptive versus unsupported.
Because the mechanism (Section~\ref{sec:mechanism}) shows the label lags, the monitor is presented as
a \emph{nowcast}, not an early-warning signal --- the recognition-lag statistics of
Section~\ref{sec:instrument} are rendered alongside every state call.%
}
\reviewnote{Adam: this section is a placeholder. Fill from the monitor build (Track C) ---
forecast-evaluation machinery (PIT, Christoffersen, CRPS/QW-CRPS, Amisano--Giacomini,
Giacomini--White) and the badge scoreboard, which doubles as the paper's honest claims ledger.}

% <<< end sections/09-monitor.tex

% >>> begin sections/10-related.tex
\section{Related literature}
\label{sec:related}

\aidraft{%
\textbf{Regime and jump-model allocation.}
The direct targets are the statistical-jump-model allocation papers: \citet{nystrup2020eswa,nystrup2021}
on the estimator, and \citet{shu2024} (single-asset downside overlay) and \citet{shu2025dynamic}
(multi-asset, expected-return-conditioned optimization), with the model-predictive-control extension
\citep{jmmpc2025}.
Our quarrel is not with the estimator --- which we adopt and validate --- but with the baselines these
papers use to establish economic value: buy-and-hold, minimum-variance, and equal-weight, none
exposure-matched or reactive.%
}

\aidraft{%
\textbf{The deflation genre.}
\citet{cederburg2020} show that volatility-managed portfolios fail out of sample across a large factor
cross-section; our result is adjacent --- a regime-conditioning strategy that dies against the same
reactive incumbents --- and inherits the same warning about a factor-timing chapter.%
}

\aidraft{%
\textbf{Reactive incumbents.}
The bars we hold the label to are the volatility-timing and volatility-managed literatures
\citep{fko2001,moreiramuir2017}, with \citet{barroso2015} the robust exception (momentum
volatility-scaling), and the time-series-momentum construction of \citet{moskowitz2012} underlying the
trend proxy of Section~\ref{sec:race4}.
The momentum-crash dynamic of \citet{danielmoskowitz2016} appears in the descriptive anatomy.%
}
\needcite{DeMiguel et al.\ 2024 if used; Kritzman--Page--Turkington 2012 (Regime Shifts).}

% <<< end sections/10-related.tex

% >>> begin sections/11-discussion.tex
\section{Discussion --- scope and what would falsify us}
\label{sec:discussion}

\aidraft{%
Our claims are deliberately bounded: daily-horizon decisions, a single developed equity market,
honest costs and next-close execution.
We make no prescriptive recommendation --- volatility targeting and EWMA estimation won \emph{matched}
races as baselines, not as endorsed strategies, and \citet{cederburg2020} shows they are themselves
fragile out of sample.
The paper deflates a claim and validates an instrument; it prescribes no strategy.%
}

\aidraft{%
The results point at where regime value could still survive: uses where lag-tolerance is
\emph{structural} --- decision cadences slow enough that a reactive estimator's speed advantage
compresses --- and information sources genuinely orthogonal to the volatility axis.
We report that the leading candidate for the latter, implied-volatility conditioning, collapses onto
the volatility axis in our data, and that a better \emph{detector} of the same return features cannot
help, because the recognition lag and the information content are both unchanged.%
}
\stub{Name the survivor hypotheses precisely (monthly/quarterly cadence; a genuinely new information
source) and state the exact control each would have to clear. Add the falsification conditions for our
own thesis.}

% <<< end sections/11-discussion.tex

% >>> begin sections/12-reproducibility.tex
\section{Reproducibility}
\label{sec:repro}

\aidraft{%
Every economic result is backed by a public git history in which the preregistration freeze commit
provably precedes the one-look run.
The frozen artifacts are the chapter-1 overlay (\texttt{results/stage1.csv},
\texttt{oos\_labels.csv}), the chapter-2 allocation battery (\texttt{results/allocation\_*}), the
calibration gate (\texttt{results/calibration\_gate.csv}), and the cross-asset descriptive battery
(reproducible at commit \texttt{d022c09}).
The estimator, walk-forward pipeline, and baselines are the same code for the synthetic and real runs;
the test suite passes, and the results dashboard regenerates every exhibit.%
}
\stub{Fill exact commit hashes / tags and the one-command reproduction steps before submission.}

% <<< end sections/12-reproducibility.tex


\bibliographystyle{plainnat}
\bibliography{refs}

\end{document}
